%% Général %%

Dès leur création, les institutions patrimoniales font face à plusieurs facteurs d'influence sur leur traitement documentaire. Nous en avons soulevé les principaux concernant l'institut national de l'audiovisuel. Certains de ces facteurs ont affecté d'autres institutions, c'est particulièrement le cas des nouvelles technologies. Au fil des années, la documentation et le traitement documentaire se sont continuellement adaptés à ses évolutions, des balbutiements de l'informatique à l'avènement de l'intelligence artificielle. D'autres facteurs sont eux plus spécifiques à l'institution. Dans ce cas, la fragilité des supports conservés par l'INA ont motivé une campagne de numérisation exceptionnelle. D'un point de vue plus institutionnel, quand l'institut national de l'audiovisuel se voit confier le dépôt légal des diffusions télévisées, le traitement documentaire s'en voit affecter de diverses manières. De façon directe quand il faut adapter la description des documents à de nouveaux publics et de nouveaux usages mais aussi de façon indirecte. Quand il est affecté par la séparation des services patrimoniaux de l'institution pour séparer les archives professionnelles des diffusions du dépôt légal puis quand le budget impose une retenue dans la sélection des documents traités et le traitement qui leur consacré. 

%% INA - Première partie %%

Ainsi, si nous tentons de retracer les principales politiques documentaires et leurs perturbations nous commençons à la première cinémathèque de l'ORTF. La politique documentaire est à des fins commerciales, les documents sont traités en stocks et les informations que l'on y trouve servent les cinémathécaire de l'ORTF pour retrouver facilement des images d'illustrations pour les journaux télévisés. Les fiches correspondant à cette politique crées avant 1975 ont été reprises par des prestataires et vérifiées par des documentalistes dans les années 90. A partir de 1975, les premières bases de données font leur apparition. Dans un premier temps suivant la logique de bibliothèque puis s'adaptant aux nouvelles normes. Entre les deux bases de données, lors du versement, certains champs ne correspondant pas ont été versés de la première base à la deuxième. En 1992 la loi du dépôt légal est voté, la conservation et la communication patrimoniale à des fins scientifique s'inscrit dans les missions de l'INA. Tout en continuant sa conservation et communication à des fins commerciales, l'INA met en place une nouvelle politique documentaire appliquée au flux. Cette nouvelle politique est au service de ce nouveau public de chercheurs qui étudient autant les émissions que leur diffusion. Pour cela les professionnels de l'INA se séparent en deux services et utilisent outils de production et de consultations, appuyés sur différentes bases de données. Pour faire face à l’afflux massif de données, les professionnels de l'INA mettent en place différents types de traitements. C'est par ce prisme que le traitements documentaire est désormais affecté car une fois mis en place, ils ne changeront presque plus. Ce qui changera, c'est la question de programme, quel chaîne se voit assigner quel traitement. Pour répondre à cette question, de nombreux critères rentrent en compte : quelle quantité de données l'INA collecte-t-elle, combien en conserve-t-elle, quel temps son budget lui permet-il d'assigner à leur traitement. En 2020, les deux services patrimoniaux fusionnent mais les deux politiques documentaires persistent. Malgré une communication facilitées entre les responsables de fonds dévolus et non dévolus, les outils de production et de consultation restent résolument doubles. Et ce alors que le traitement documentaire continu d'évoluer face à l'intégration de l'IA dans les pratiques documentaires. 

%% Reprise d'antériorité %%

Et pour conclure, parmi toutes ces notices issues de politiques documentaires différentes, affectées par plusieurs facteurs internes ou externes à l'institution, certaines de ces notices anciennes ont potentiellement été l’œuvre de reprises d'antériorité, elles aussi effectuées, selon les politiques documentaire et les situations institutionnelles, technologiques et budgétaires en vigueur à leur époque, éliminant l'indexation utilisée à la création de la première, surtout pour les fiches papiers postérieures à 1975 dont nous n'avons plus de traces, au contraires des fiches papier précédent 1975. Tout cela alors que très peu de ces changements ne sont documentés. 

%% Technologies - Les solutions déjà mises en place %%

De plus, le traitement documentaire en vigueur n'est pas figé dans le marbre, il est en constante évolution, notamment avec l'utilisation de l'intelligence artificielle générative  qui permet d'améliorer l'accessibilité aux collections tout en remettant en question le métier de documentaliste. Il faut prévoir son impact sur le traitement documentaire.   

%% Ouverture à la partie qui suit %%

C'est avec toutes ces informations en tête que nous avons élaboré un outil ayant pour objectif de faciliter l'accessibilité aux collections de l'INA. Désormais nous avons conscience des enjeux que représentent l’accessibilité aux collections dans les institutions patrimoniales et plus particulièrement à l'INA ainsi que l'origine des de ces enjeux. Nous disposons de toutes les informations nécessaire pour construire un outil pouvant mettre à contribution les outils facilitant l'accès aux collections tout en offrant une solution pour tous les défauts résultant de l'évolution du traitement documentaire. 